\makeatletter \def\input@path{{../../}} \makeatother
\documentclass{article}
\usepackage[utf8]{inputenc}
\usepackage{amsmath, amsthm, amssymb, mathpazo, isomath, mathtools}
\usepackage{subcaption,graphicx,pgfplots}
\usepackage{fullpage}
\usepackage{booktabs}
\usepackage{hyperref}
\usepackage{algorithm, algorithmic}
\usepackage{mathtools}
\usepackage{todonotes}
\usepackage{tcolorbox}
\usepackage{totcount,xfp}
\newtotcounter{totalpts}
%\setcounter{totalpts}{0} % Contador de puntos en homework

\newcommand{\getTotal}{\number\totvalue{totalpts} }
\newcommand{\instruccion}{\begin{tcolorbox}[colback=blue!5!white,colframe=blue!75!black!70!white,title=Instructions]
The homework must be developed by one or two people in a .pdf report submitted through the designated drop box on the Canvas platform, where you must also show the developed code if applicable. For your convenience, you may submit the homeworks in a Jupyter Notebook, so it is more comfortable to show the code. In any case, a single file must be submitted as the answer to the homework. The late policy will be: a linear factor will be calculated that is 1 at the time of delivery and 0 48 hours later. This will multiply your obtained score.


The homework has \getTotal points. In each question, the given score will be: the total if it is correct, half if it has a conclusive development but contains errors, and 0.0 points if it has little to no development.

The use of language models (like ChatGPT) is allowed to support the development of the homework, although we suggest trying each question a bit first to generate doubts. While we cannot control if you ask the model directly for the answer, we believe it is better for your learning process to ask for suggestions rather than solutions. For this reason, if you decide to use it, you must attach to your homework another pdf showing the prompts used. We suggest using prompts of this style: "I am a student of [Major] and I am developing the following exercise for a homework in the course [Course]: [Copy exercise]
  And we have seen so far the contents in the attached slides [attach pdfs of the lectures]. Given this content, I would like to ask you for a suggestion to solve [explain difficulty], since I am confused about [write doubt]. As an answer, I do not want you to give me the solution, but a guide to help me start developing my answer."

  If you have doubts and/or comments, please contact us by email or Canvas. It is possible that the formulation has errors.
\end{tcolorbox}}


\newcommand{\code}[1]{\texttt{#1}}
% Logic: To create a counter with the points in each question and accumulate that in a total variable. Now I simply set the full threshold at the beginning to 75% and it gives the computed value automatically.
\newcommand{\pts}[1]{\if#11\textbf{[#1 point]}\else\textbf{[#1 points]}\fi\addtocounter{totalpts}{#1}}


\title{Homework 2}
%\author{Nicol\'as A Barnafi}
\date{}

\renewcommand{\vec}{\vectorsym}
\newcommand{\mat}{\matrixsym}
\newcommand{\ten}{\tensorsym}
\DeclareMathOperator{\grad}{\nabla}
\DeclareMathOperator{\dive}{\text{div}}
\DeclareMathOperator{\curl}{\text{curl}}
\DeclareMathOperator{\tr}{\text{tr}}
\DeclareMathOperator{\sym}{\text{sym}}
\newtheorem{remark}{Remark}
\newtheorem{definition}{Definition}
\newcommand{\R}{\mathbb{R}}
\newcommand{\D}{\mathcal{D}}

\newcommand{\tin}{\text{in}}
\newcommand{\ton}{\text{on}}

\newtheorem{theorem}{Theorem}
\newtheorem{lemma}{Lemma}

\begin{document}

\maketitle
 
\instruccion

\begin{enumerate}
    \item Consider the bistable problem with diffusion in time and space $\Omega_T = (0,T)\times (0,1)$:
            $$ \dot u - \mu \partial_{xx} u - u(1-u^2) = f. $$
            As data and boundary conditions, consider:
            \begin{itemize}
                \item $f(t,x) = I_{x\in (0.4,0.6)}I_{t<1}$, where $I$ is an indicator function.
                \item An initial condition given by $u(0,x) = 0$ for all $x$ in $(0,1)$. 
                \item Homogeneous Neumann boundary conditions for all $t>0$. 
            \end{itemize}
            For this model:
            \begin{enumerate}
                \item\pts{1} Propose three time discretization formulations: (i) explicit, (ii) implicit, and (iii) semi-implicit. 
                \item\pts{1} Propose a space and time discretization scheme, and prove that it is consistent.
                \item\pts{1} Formulate an iterative method (Newton or fixed point) to solve the implicit formulation and comment on the convergence properties of the proposed method. 
                \item\pts{1} The method proposed in the previous question requires solving a linear system in each iteration. Using von Neumann analysis, propose hypotheses for the discrete solution and the discretization parameters that allow obtaining a stable approximation in each iteration.
                \item\pts{2} Implement the proposed semi-implicit method and plot the evolution of the solution for $\mu=0.1$ over a time horizon of $T=5$. Comment on your results.
                \item\pts{3} Implement the proposed implicit method and plot the evolution of the solution for $\mu=0.1$. Comment on your results and compare them with those obtained in the previous point.
            \end{enumerate}

    \item In this question, we will play with the definition of distributions. We will use the notation $\mathcal D(\R) = C_0^\infty(\R)$, the space of smooth functions with compact support in $\R$.
        \begin{enumerate}
            \item\pts{1} We define the injection of integrable functions into distributions $T:L^1(\R) \to \left(\mathcal D(\R)\right)'$ according to the action
                $$ (Tf)(\varphi)  = \left\langle f, \varphi\right\rangle_{\mathcal D'\times \mathcal D}.$$
                Prove that this injection is continuous. 
            \item\pts{1} Prove that, given $x_0\in \R$, the sequence induced by the function $I_n(x) = nI_{(x_0-1/2n, x_0+1/2n)}$, i.e. $\{I_n\}_n$, converges as a distribution to the Dirac delta $\delta_{x_0}$. 
            \item\pts{1} Prove that there is no function $f$ in $L^1$ such that $T_f = \delta_0$. 
            \item\pts{2} We know that given a vector function $\vec F:\Omega \to \R^d$, its divergence is given by $\dive{\vec F}=F_{i,i}$. To extend this notion to distributions, prove that given a domain $\Omega$, the spaces $(\mathcal D(\Omega, \R^d))'$ and $[(\mathcal D(\Omega))']^d$ are homeomorphic\footnote{$X^3 = X\times X\times X$} and with this, construct a definition of divergence in $(\mathcal D(\Omega, \R^d))'$. Hint: $\int_\Omega (\dive F)\varphi = \int_{\partial\Omega} \varphi F\cdot \vec n - \int_\Omega F\cdot \grad \varphi$. 
            \item\pts{1} Extend the construction from the previous question to define a distributional $\curl$ for the case $d=3$. 
            \item\pts{2} Consider the function $\Phi:\R^3\to \mathbb C$ given by
                    $$ \Phi(x) = \frac{1}{4\pi|x|} e^{-ik |x|}. $$
                Show that the following equality holds in $\mathcal D'(\R^3)$:
                        $$ -\Delta \Phi - k^2 \Phi = \delta_0, $$
                i.e. in the sense of distributions, where $\delta_0$ is the Dirac delta at $x=0$. 
        \end{enumerate}
\end{enumerate}

\todo[inline,color=white!90!black]{\textbf{Note: } We will open a forum on Canvas to review any typos and/or errors in the formulation.}
\end{document}
